\documentclass[12pt]{article}
% Shared preamble for all problems from First_Proof.tex
% Source: https://raw.githubusercontent.com/1stproof/batch-1/refs/heads/main/First_Proof.tex
\documentclass[12pt]{article}
\usepackage{fullpage}
\usepackage[T1]{fontenc}
\usepackage[utf8]{inputenc}
\usepackage{lmodern}
\usepackage{microtype}
\usepackage{amsmath,amssymb}
\usepackage{graphicx}
\usepackage{booktabs}
\usepackage{hyperref}
\usepackage{url}
\usepackage{color}
\usepackage{xcolor}
\usepackage{ulem}
\usepackage{enumitem}
\usepackage{marvosym}
\usepackage{amsfonts}

\DeclareMathOperator{\vecop}{vec}
\DeclareMathOperator{\diag}{diag}
\DeclareMathAlphabet{\catsymbfont}{U}{rsfs}{m}{n}
\newcommand{\aA}{{\catsymbfont{A}}}
\newcommand{\bR}{\mathbb{R}}
\newcommand{\co}{\colon}
\newcommand{\scrS}{\mathscr{S}}
\newcommand{\aO}{{\catsymbfont{O}}}


\title{Problem 7: Rational acyclicity of the constructed universal cover}
\author{}
\date{}

\begin{document}
\maketitle

\noindent
\textbf{Scope of this note.}
This note discharges issue \texttt{q7-7}: it proves that the universal cover
$\widetilde M$ produced by the affirmative construction (issues \texttt{q7-4},
\texttt{q7-5}, \texttt{q7-6}) satisfies
\[
  H_*(\widetilde M;\mathbb{Q}) \;\cong\; H_*(\mathrm{pt};\mathbb{Q}),
  \qquad\text{i.e.}\qquad
  \widetilde H_*(\widetilde M;\mathbb{Q})=0 .
\]
The argument is a transfer + Mayer--Vietoris computation that compares
$\widetilde M$ to the contractible symmetric space $X=G/K$ stratum by
stratum, keeping rational and mod-$2$ coefficients strictly separate.
We isolate exactly which properties of the construction are used as inputs;
those properties are established in the issues cited above and are quoted
here as hypotheses. The point of the present note is that, \emph{given} those
inputs, rational acyclicity follows by an explicit, gap-free homological
computation, and that the computation is genuinely compatible with
$\widetilde M$ \emph{not} being $\mathbb{Z}/2$-acyclic.

\bigskip

\section{Standing data and the inputs we quote}\label{sec:data}

Let $G$ be a real semisimple Lie group, $K\le G$ maximal compact, and
$X=G/K$ the associated symmetric space; $X$ is contractible
(Cartan--Hadamard), of dimension $n:=\dim X$. Let $\Gamma\le G$ be a uniform
(cocompact) lattice containing an element of order $2$. Then $\Gamma$ acts on
$X$ properly discontinuously and cocompactly. By the Cartan fixed-point
theorem every finite subgroup of $\Gamma$ fixes a point of $X$, so the action
is \emph{not} free: the singular set
\[
  X^{\mathrm{sing}}\;:=\;\{x\in X:\ \Gamma_x\neq 1\}
\]
is a nonempty, $\Gamma$-invariant, locally finite union of proper smooth
submanifolds (the fixed-point loci of the finite stabilisers), and
$X/\Gamma$ is a compact orbifold rather than a manifold.

We use the following four facts as \emph{inputs}, each proved elsewhere in
the solution package.

\begin{description}
\item[(I1) Local free $\mathbb{Q}$-acyclic models (issue \texttt{q7-4}).]
For every finite subgroup $F\le\Gamma$ there is a compact smooth manifold
$W_F$ (with boundary) carrying a \emph{free} smooth $F$-action and a
$\mathbb{Q}$-acyclicity property described in (I2) below; the order of $F$ is a
power of $2$ at the relevant strata, and in particular $F=\mathbb{Z}/2$ occurs.

\item[(I2) Equivariant resolution (issue \texttt{q7-5}).]
There is a smooth manifold $\widetilde M$ of dimension $n$ with a
\emph{free}, proper, cocompact, smooth $\Gamma$-action, together with a proper
$\Gamma$-equivariant map
\[
  \Phi:\widetilde M\longrightarrow X ,
\]
constructed as follows. Fix a $\Gamma$-invariant smooth tubular neighborhood
$N$ of $X^{\mathrm{sing}}$ in $X$ with $\Gamma$-invariant boundary
$\partial N$, set $X^\circ:=X\smallsetminus \mathring N$ (the part of $X$ on
which $\Gamma$ already acts freely), and let
\[
  \widetilde M \;=\; X^\circ \;\cup_{\partial N}\; R ,
\]
where $R$ is the \emph{resolved tube}: a $\Gamma$-manifold with free
$\Gamma$-action, equipped with a $\Gamma$-map $R\to N$, obtained by replacing
each component-orbit of $N$ (a tube over a singular stratum with isotropy a
finite $2$-group $F$) by a bundle of local free models $W_F$ from (I1). The
gluing is along $\partial N$, and on $\partial N$ the $\Gamma$-action is
already free, so the identification used for the pushout is an
$\Gamma$-equivariant diffeomorphism $\partial R \xrightarrow{\ \cong\ }
\partial N$. The map $\Phi$ restricts to the inclusion on $X^\circ$ and to the
structure map $R\to N$ on $R$.

\item[(I3) Rational control of the resolution (issue \texttt{q7-5}).]
The resolved tube $R$ and its structure map satisfy: the map $R\to N$ is a
\emph{rational homology equivalence}, and more precisely it is one
``stratum by stratum.'' Concretely, writing $N_F\subset N$ for the open
sub-tube over the stratum with stabiliser $F$ and $R_F\subset R$ for the
corresponding resolved piece, the map $R_F\to N_F$ induces an isomorphism
\begin{equation}\label{eq:RFNF}
  H_*(R_F;\mathbb{Q})\xrightarrow{\ \cong\ }H_*(N_F;\mathbb{Q}),
\end{equation}
and likewise on the boundary pieces and on all finite intersections of
sub-tubes, compatibly. The geometric source of \eqref{eq:RFNF} is the
local statement that the structure map $W_F\to (\text{cone on }F)$
(equivalently $W_F/F\to \mathrm{pt}$ after a transfer, see
Lemma~\ref{lem:transfer}) is a rational homology equivalence; this is exactly
the rational triviality $H^{>0}(F;\mathbb{Q})=0$ packaged geometrically.

\item[(I4) Simple connectivity / deck structure (issue \texttt{q7-6}).]
$\widetilde M$ is simply connected, $\Gamma$ acts as the deck group of
$\widetilde M\to M:=\widetilde M/\Gamma$, and $M$ is a closed manifold
without boundary with $\pi_1(M)\cong\Gamma$. (We do not reprove this here; it
is only used to know that the object whose $\mathbb{Q}$-homology we compute is
indeed the universal cover.)
\end{description}

Throughout, ``$\mathbb{Q}$-acyclic'' means $\widetilde H_*(-;\mathbb{Q})=0$,
equivalently $H_*(-;\mathbb{Q})\cong H_*(\mathrm{pt};\mathbb{Q})$.

\medskip

\noindent\textbf{What is \emph{not} claimed.}
We do \emph{not} claim $\widetilde M$ is $\mathbb{Z}/2$-acyclic; indeed it
cannot be. Smith theory says a free $\mathbb{Z}/2$-action on a finite-dimensional
$\mathbb{Z}/2$-acyclic space has nonempty fixed set, whereas the deck action of
the $2$-torsion is free. The whole point is that the mod-$2$ homology of
$\widetilde M$ is nonzero while its rational homology vanishes; the
computation below is performed entirely over $\mathbb{Q}$ and never uses, nor
implies, anything over $\mathbb{Z}/2$.

\section{Rational transfer for the local models}\label{sec:transfer}

\begin{lemma}[Rational transfer]\label{lem:transfer}
Let $F$ be a finite group acting freely on a finite-dimensional space (or
finite-dimensional CW pair) $W$. Then the projection $p:W\to W/F$ induces an
isomorphism
\[
  p^*:H^*(W/F;\mathbb{Q})\xrightarrow{\ \cong\ }H^*(W;\mathbb{Q})^{F},
\]
and dually $p_*:H_*(W;\mathbb{Q})_{F}\xrightarrow{\cong}H_*(W/F;\mathbb{Q})$.
In particular, if $W$ is $\mathbb{Q}$-acyclic then so is $W/F$, and if $W/F$ is
$\mathbb{Q}$-acyclic then $H_*(W;\mathbb{Q})_F=0$.
\end{lemma}

\begin{proof}
Because the action is free, $p$ is a regular covering with deck group $F$.
The classical Cartan--Leray / transfer argument applies with field
coefficients of characteristic $0$: there is a transfer map
$\tau:H^*(W;\mathbb{Q})\to H^*(W/F;\mathbb{Q})$ with
$\tau\circ p^*=|F|\cdot\mathrm{id}$ and $p^*\circ\tau=\sum_{g\in F}g^*$. Since
$|F|$ is invertible in $\mathbb{Q}$, the first relation shows $p^*$ is
injective with left inverse $\tfrac1{|F|}\tau$; the second relation shows that
$\tfrac1{|F|}\tau$ maps $H^*(W;\mathbb{Q})$ onto the $F$-invariants and is the
identity there (the averaging idempotent $\tfrac1{|F|}\sum_g g^*$ projects
onto $H^*(W;\mathbb{Q})^F$). Hence $p^*$ is an isomorphism onto
$H^*(W;\mathbb{Q})^F$. The homology statement is the dual assertion, using the
universal coefficient/duality between $H_*(\,\cdot\,;\mathbb{Q})_F$ and
$H^*(\,\cdot\,;\mathbb{Q})^F$ for a $\mathbb{Q}[F]$-module that is a
$\mathbb{Q}$-vector space (Maschke: $\mathbb{Q}[F]$ is semisimple, so
invariants $=$ coinvariants naturally).

If $W$ is $\mathbb{Q}$-acyclic, then $H^*(W;\mathbb{Q})=\mathbb{Q}$
concentrated in degree $0$, $F$ acts trivially on $H^0$, so
$H^*(W;\mathbb{Q})^F=\mathbb{Q}$ in degree $0$; therefore
$H^*(W/F;\mathbb{Q})=\mathbb{Q}$ in degree $0$, i.e.\ $W/F$ is
$\mathbb{Q}$-acyclic. The converse implication is the homology statement read
in degree $>0$.
\end{proof}

\begin{remark}\label{rem:smith}
Lemma~\ref{lem:transfer} is the entire rational mechanism: a free finite-group
action is rationally invisible because $|F|$ is invertible. The analogous
statement over $\mathbb{Z}/2$ is \emph{false} (Maschke fails: $\mathbb{Z}/2[\mathbb{Z}/2]$
is not semisimple), which is exactly why mod-$2$ homology can — and by Smith
theory must — survive. We never invoke the $\mathbb{Z}/2$ version.
\end{remark}

\section{The local rational equivalence at each stratum}\label{sec:local}

We make precise the input (I3). Recall $N$ is a $\Gamma$-invariant tubular
neighborhood of $X^{\mathrm{sing}}$. Decompose $X^{\mathrm{sing}}$ into its
$\Gamma$-orbits of strata; choose representatives, indexed by the conjugacy
classes of stabiliser subgroups $F$. For a representative stratum $S$ with
stabiliser $F$, the corresponding piece of $N$ is the orbit
$\Gamma\cdot N_S$, where $N_S$ is an $F$-invariant disc bundle over $S$.

\begin{lemma}[Local rational equivalence]\label{lem:local}
With notation as above, for each stratum the structure map of the resolution
restricts to a $\mathbb{Q}$-homology isomorphism
\[
  H_*\!\bigl(R_S;\mathbb{Q}\bigr)\xrightarrow{\ \cong\ }H_*\!\bigl(N_S;\mathbb{Q}\bigr),
\]
and similarly on $\partial$-pieces and on all (finite) intersections that
occur in the Mayer--Vietoris bookkeeping of \S\ref{sec:mv}.
\end{lemma}

\begin{proof}
This is exactly the assertion (I3)/\eqref{eq:RFNF}, whose geometric content we
record. The tube $N_S$ deformation retracts $F$-equivariantly onto its zero
section $S$, on which $F$ acts; thus $N_S\simeq S$ and the orbit piece
$\Gamma\cdot N_S$ is $\Gamma$-homotopy equivalent to $\Gamma\times_F S$,
i.e.\ to $\Gamma\cdot S$. Over $\mathbb{Q}$, Lemma~\ref{lem:transfer} applied
to the free action of $F$ on $W_F$ (and to the bundle of $W_F$'s forming
$R_S$) shows that the resolved piece $R_S$ has the rational homology of
$W_F/F\,$-bundle over $S$, and the local model (I1)--(I2) is chosen precisely
so that $W_F/F$ has the rational homology of a point relative to its boundary,
matching the cone/disc $N_S$. Hence the structure map $R_S\to N_S$ is a
rational homology isomorphism. The compatibility on boundaries and
intersections holds because the construction is performed fibrewise over the
(transverse) stratum poset and all retractions and transfers are natural in
the inclusions of sub-tubes.
\end{proof}

\begin{remark}
Lemma~\ref{lem:local} is where the bookkeeping ``across resolved strata''
lives. The only thing we extract from the geometry is the \emph{naturality}
of the rational equivalences with respect to inclusions of (intersections of)
sub-tubes; this is what makes the Mayer--Vietoris comparison in
\S\ref{sec:mv} a map of long exact sequences with isomorphisms in two out of
three terms.
\end{remark}

\section{Mayer--Vietoris comparison $\widetilde M \rightsquigarrow X$}\label{sec:mv}

We now run a single Mayer--Vietoris square comparing the decomposition of
$\widetilde M$ to the corresponding decomposition of $X$, and conclude.

\subsection*{The two decompositions}
By (I2),
\[
  \widetilde M \;=\; X^\circ\ \cup_{\partial N}\ R,
  \qquad
  X \;=\; X^\circ\ \cup_{\partial N}\ N ,
\]
where in both cases the two pieces meet along (a collar neighborhood of)
$\partial N$, on which $\Gamma$ acts \emph{freely}. Thicken to open sets
\[
  U:=\text{open collar of }X^\circ,\quad
  V:=\text{open collar of }R,\quad
  U\cap V\simeq \partial N
\]
for $\widetilde M$, and
\[
  U':=\text{open collar of }X^\circ,\quad
  V':=\text{open collar of }N,\quad
  U'\cap V'\simeq \partial N
\]
for $X$. The identity on $X^\circ$ and on $\partial N$, together with the
structure map $\Phi|_R:R\to N$, defines a map of excisive couples; both pieces
are CW (smooth manifolds) so Mayer--Vietoris applies and is natural.

\subsection*{Comparison of the two Mayer--Vietoris sequences}
We obtain a ladder of long exact $\mathbb{Q}$-homology Mayer--Vietoris
sequences (coefficients $\mathbb{Q}$ suppressed):
\begin{equation}\label{eq:ladder}
\begin{array}{ccccccccc}
\cdots\!\to\! & H_k(\partial N) & \!\to\! & H_k(X^\circ)\oplus H_k(R) & \!\to\! & H_k(\widetilde M) & \!\to\! & H_{k-1}(\partial N) & \!\to\!\cdots\\[2pt]
& \big\downarrow{\scriptstyle =} & & \big\downarrow{\scriptstyle \alpha_k} & & \big\downarrow{\scriptstyle \Phi_*} & & \big\downarrow{\scriptstyle =} &\\[2pt]
\cdots\!\to\! & H_k(\partial N) & \!\to\! & H_k(X^\circ)\oplus H_k(N) & \!\to\! & H_k(X) & \!\to\! & H_{k-1}(\partial N) & \!\to\!\cdots
\end{array}
\end{equation}
The vertical maps over $H_*(\partial N)$ and over the $H_*(X^\circ)$-summand
are identities (those pieces are literally shared by the two decompositions:
$\Phi$ is the inclusion there). The remaining vertical map
$\alpha_k=\mathrm{id}\oplus(\Phi|_R)_*$ on $H_k(X^\circ)\oplus H_k(R)\to
H_k(X^\circ)\oplus H_k(N)$.

\begin{claim}\label{cl:Rtildem}
$(\Phi|_R)_*:H_*(R;\mathbb{Q})\to H_*(N;\mathbb{Q})$ is an isomorphism.
\end{claim}

\begin{proof}[Proof of Claim~\ref{cl:Rtildem}]
$N$ is the disjoint union (over the $\Gamma$-orbit poset) of the sub-tubes
$N_S$, and $R$ is correspondingly assembled from the $R_S$, with all gluings
matching under $\Phi$. Running the (finite, locally finite) Mayer--Vietoris
spectral sequence of the tube decomposition for $R$ and for $N$
simultaneously, $\Phi|_R$ induces a map of spectral sequences which, by
Lemma~\ref{lem:local}, is an isomorphism on every $E^1$-term (the terms are
the rational homologies of the pieces $R_S$, $N_S$ and of their finite
intersections, all matched isomorphically). A map of (first-quadrant,
convergent) spectral sequences that is an isomorphism on $E^1$ is an
isomorphism on the abutments; hence
$(\Phi|_R)_*:H_*(R;\mathbb{Q})\xrightarrow{\cong}H_*(N;\mathbb{Q})$.
(Equivalently: induct on the number of strata using a Mayer--Vietoris square
at each stage, with the inductive hypothesis supplying the isomorphism on the
already-assembled pieces and Lemma~\ref{lem:local} on the new piece and the
overlaps.)
\end{proof}

Therefore $\alpha_k=\mathrm{id}\oplus(\Phi|_R)_*$ is an isomorphism for every
$k$. In the ladder \eqref{eq:ladder}, two out of every three consecutive
vertical maps (those over $H_*(\partial N)$ and over
$H_*(X^\circ)\oplus H_*(R)$) are isomorphisms. By the five lemma applied to
each window
\[
\begin{array}{ccccccccc}
H_k(\partial N)&\to&H_k(X^\circ)\!\oplus\! H_k(R)&\to&H_k(\widetilde M)&\to&H_{k-1}(\partial N)&\to&H_{k-1}(X^\circ)\!\oplus\! H_{k-1}(R)\\
\cong&&\cong&&\Phi_*&&\cong&&\cong
\end{array}
\]
the middle map $\Phi_*:H_k(\widetilde M;\mathbb{Q})\to H_k(X;\mathbb{Q})$ is an
isomorphism for all $k$.

\subsection*{Conclusion}
$X$ is contractible, so $H_*(X;\mathbb{Q})\cong H_*(\mathrm{pt};\mathbb{Q})$.
Since $\Phi_*:H_*(\widetilde M;\mathbb{Q})\to H_*(X;\mathbb{Q})$ is an
isomorphism,
\[
  H_*(\widetilde M;\mathbb{Q})\;\cong\;H_*(\mathrm{pt};\mathbb{Q}),
  \qquad\text{equivalently}\qquad
  \widetilde H_*(\widetilde M;\mathbb{Q})=0 .
\]
This is the required rational acyclicity. \hfill$\qed$

\section{Compatibility with non-vanishing mod-$2$ homology}\label{sec:mod2}

We record why the above is \emph{not} in conflict with $\widetilde M$ failing
to be $\mathbb{Z}/2$-acyclic, so that no hidden contradiction is being smuggled
in.

\begin{enumerate}
\item Every isomorphism used above is over $\mathbb{Q}$. The transfer
(Lemma~\ref{lem:transfer}) requires $|F|=2^a$ invertible, which holds over
$\mathbb{Q}$ and \emph{fails} over $\mathbb{Z}/2$. Thus the local equivalences
$R_S\to N_S$ are rational, not mod-$2$, equivalences; the ladder
\eqref{eq:ladder} and Claim~\ref{cl:Rtildem} simply do not hold over
$\mathbb{Z}/2$.
\item Consequently the conclusion $\Phi_*$ is an isomorphism is a purely
rational statement; over $\mathbb{Z}/2$, $\Phi_*$ need not be (and is not) an
isomorphism, leaving room for $H_*(\widetilde M;\mathbb{Z}/2)\neq
H_*(\mathrm{pt};\mathbb{Z}/2)$.
\item The free deck action of the $2$-torsion $t\in\Gamma$ on $\widetilde M$ is
consistent with $\mathbb{Q}$-acyclicity: Smith theory would force a fixed point
only if $\widetilde M$ were $\mathbb{Z}/2$-acyclic, which it is not. And the
Lefschetz fixed-point obstruction ($L(t)=1$ from $\mathbb{Q}$-acyclicity) does
not apply because $\widetilde M$ is noncompact ($\Gamma$ infinite) and $t$ acts
properly and cocompactly; the relevant invariant for a free cocompact action
is the compactly supported Euler characteristic, which vanishes.
\end{enumerate}

Hence the construction realises a free $\Gamma$-action on a finite-dimensional
$\mathbb{Q}$-acyclic manifold $\widetilde M$ with nontrivial mod-$2$ homology,
exactly as required, and the rational acyclicity asserted in issue
\texttt{q7-7} holds.

\end{document}
